% Bagian: first-order-logic
% Bab: proof-systems
% Subbagian: axiomatic-deduction

\documentclass[../../../include/open-logic-section]{subfiles}

\begin{document}

\iftag{FOL}
      {\olfileid[id]{fol}{prf}{axd}}
      {\olfileid[id]{pl}{prf}{axd}}

\olsection{\usetoken{P}{derivation} Aksiomatik}

Sistem !!{derivation}s aksiomatik merupakan sistem !!{derivation} logis
yang paling tua dan paling sederhana. Di dalamnya, !!{derivation}s
hanyalah barisan !!{sentence}s. Suatu barisan !!{sentence}s dianggap
sebagai !!{derivation} yang benar jika setiap !!{sentence}~$!A$ di
dalamnya memenuhi salah satu syarat berikut:
\begin{enumerate}
\item $!A$ merupakan aksioma, atau
\item $!A$ merupakan !!a{element} dari suatu himpunan~$\Gamma$ yang
  beranggotakan !!{sentence}s, atau
\item $!A$ dijustifikasi oleh aturan inferensi.
\end{enumerate}
Untuk menjadi aksioma, $!A$ harus memiliki bentuk salah satu dari
sejumlah skema !!{sentence} yang tetap. Terdapat banyak himpunan skema
aksioma yang menghasilkan sistem !!{derivation} logika orde pertama
yang memadai (sahih dan lengkap). Beberapa di antaranya disusun
menurut penghubung yang diaturnya, misalnya skema
\[
!A \lif (!B \lif !A) \qquad !B \lif (!B \lor !C) \qquad (!B \land !C) \lif !B
\]
merupakan aksioma umum yang mengatur $\lif$, $\lor$, dan~$\land$.
Beberapa sistem aksioma mengupayakan jumlah aksioma yang minimal.
Bergantung pada penghubung yang diambil sebagai penghubung primitif,
bahkan dimungkinkan untuk menemukan sistem aksioma yang hanya terdiri
atas satu aksioma.

Aturan inferensi adalah pernyataan bersyarat yang memberikan syarat
cukup agar !!a{sentence} dalam !!a{derivation} dapat dijustifikasi.
Modus ponens merupakan salah satu aturan semacam itu yang sangat umum:
aturan ini menyatakan bahwa jika $!A$ dan $!A \lif !B$ sudah
dijustifikasi, maka $!B$ dijustifikasi. Artinya, suatu baris dalam
!!a{derivation} yang memuat !!{sentence}~$!B$ dijustifikasi, asalkan
baik $!A$ maupun $!A \lif !B$ (untuk suatu !!{sentence}~$!A$) muncul
dalam !!{derivation} sebelum~$!B$.

Relasi $\Proves$ berdasarkan !!{derivation}s aksiomatik didefinisikan
sebagai berikut: $\Gamma \Proves !A$ jika dan hanya jika terdapat
!!a{derivation} dengan !!{sentence}~$!A$ sebagai formula terakhirnya
(dan himpunan berhingga !!{sentence}s dalam !!{derivation} tersebut
yang dijustifikasi oleh~(2) di atas merupakan himpunan bagian
dari~$\Gamma$). $!A$
merupakan teorema jika~$!A$ memiliki !!a{derivation} dengan~$\Gamma$
kosong, yaitu setiap !!{sentence} dalam !!{derivation} dijustifikasi
oleh (1) atau~(3). Sebagai contoh, berikut adalah !!a{derivation} yang
menunjukkan bahwa $\Proves !A  \lif (!B \lif (!B \lor !A))$:
\begin{derivation}
  1. & $!B \lif (!B \lor !A)$ \\
  2. & $(!B \lif (!B \lor !A)) \lif (!A  \lif (!B \lif (!B \lor !A)))$\\
  3. & $!A  \lif (!B \lif (!B \lor !A))$
\end{derivation}
Pada baris~1, !!{sentence} tersebut memiliki bentuk aksioma $!A \lif
(!A \lor !B)$ (dengan peran $!A$ dan $!B$ dipertukarkan). Kalimat pada
baris~2 memiliki bentuk aksioma $!A \lif (!B \lif !A)$. Dengan
demikian, kedua baris itu dijustifikasi. Baris~3 dijustifikasi oleh
modus ponens: jika kita menyingkatnya sebagai $!D$, maka baris~2
berbentuk $!C \lif !D$, dengan $!C$ adalah $!B \lif (!B \lor !A)$,
yaitu baris~1.

Himpunan $\Gamma$ tidak konsisten jika $\Gamma \Proves \lfalse$.
Sistem aksioma lengkap juga akan membuktikan bahwa $\lfalse \lif !A$
untuk setiap~$!A$, sehingga jika $\Gamma$ tidak konsisten, maka
$\Gamma \Proves !A$ untuk setiap~$!A$.

Sistem !!{derivation}s aksiomatik untuk logika pertama kali diberikan
oleh Gottlob Frege dalam \emph{Begriffsschrift} karyanya pada tahun
1879; karena alasan ini, karya tersebut sering dianggap sebagai karya
pertama dalam logika modern. Sistem-sistem itu disempurnakan dalam
\emph{Principia Mathematica} karya Alfred North Whitehead dan Bertrand
Russell, serta oleh David Hilbert dan para mahasiswanya pada tahun
1920-an. Oleh karena itu, sistem-sistem tersebut sering disebut
``sistem Frege'' atau ``sistem Hilbert.'' Sistem-sistem itu sangat
luwes dalam arti bahwa sistem aksiomatik bagi suatu logika sering
mudah ditemukan. Karena !!{derivation}s memiliki struktur yang sangat
sederhana dan hanya satu atau dua aturan inferensi, membuktikan berbagai
hal \emph{tentangnya} juga relatif mudah. Namun, sistem-sistem itu sangat
sulit digunakan dalam praktik, yaitu sulit menemukan dan menuliskan
pembuktian.

\end{document}
